\begingroup
\let\clearpage\relax
\let\cleardoublepage\relax
\chapter*{How to read Part IV}
\endgroup
\addcontentsline{toc}{chapter}{How to read Part IV}
\label{ch:results-map}

\section*{The argument}

The two results chapters follow a feedback loop between solver and physics, and it is worth
having the shape of it before the details.

We first establish that the trained variational states reproduce every available external
energy benchmark across the regime in which such a benchmark exists. That licenses using
them as instruments rather than defending them as answers. The states are then read for the
physics they contain, and the crossover from a compact quantum dot to a finite Wigner
molecule turns out to be \emph{staged}: the cloud expands and its shells stiffen radially
well before the electrons acquire angular order within those shells.

We then turn from the state to the solver. Controlled architecture comparisons show that
message passing changes the geometry of the variational representation even where the total
energy cannot tell two networks apart: the correlator concentrates its tangent weight on
fewer and more physically interpretable directions, and the message-passing backflow avoids
a rank collapse that overtakes the conventional one. Finally we ask what governs whether
such states can be \emph{trained} at all. Architecture, optimiser and sampling measure meet
in one object---the geometry of the wavefunction's log-derivatives---and natural-gradient
preconditioning earns its cost only while useful geometric information can still be
estimated from the batch.

The loop then closes. When a generic collocation proposal loses the Wigner state, the
structural measurements of Chapter~\ref{ch:results} supply the remedy: a proposal built from
the observed shell geometry carries Markov-chain-free training to $\omega=0.0035$, where a
different limit---a transfer instability in the Slater--Jastrow balance---takes over.

\medskip
\noindent The four research questions of \S\ref{sec:research-questions} are answered in the
following places, with the energy benchmarks of \S\ref{sec:energy-overview} as the
foundation that licenses the rest.

\begin{center}\small
\begin{tabular}{@{}ll@{}}
\toprule
\textbf{RQ4 --- Physics} & \S\ref{sec:wigner-molecule} \\
\textbf{RQ1 --- Architecture} & \S\ref{sec:kernel-q1} \\
\textbf{RQ2 --- Optimiser} & \S\ref{sec:kernel-cond} \\
\textbf{RQ3 --- Markov-chain-free paradigm} & \S\ref{sec:kernel-mcmcfree}--\S\ref{sec:kernel-q3-frontier} \\
\bottomrule
\end{tabular}
\end{center}

\section*{Three training programmes, which are never compared with one another}

Three separate sets of models appear in what follows. They exist for different purposes and
are sized differently, and numbers from one are never compared with numbers from another.

\begin{center}\small
\setlength{\tabcolsep}{5pt}
\begin{tabular}{@{}p{0.19\textwidth} p{0.30\textwidth} p{0.13\textwidth} p{0.30\textwidth}@{}}
\toprule
Programme & Purpose & Size & Supports \\
\midrule
Production (SR--VMC) & best attainable state; the physics of Ch.~\ref{ch:results} &
  $137$--$236$k & energy and structural claims \\
Architecture study & controlled separable vs.\ message-passing comparison &
  matched ladder, $20$--$164$k & architecture claims \\
Collocation & Markov-chain-free training &
  $25$--$49$k & training-paradigm claims \\
\bottomrule
\end{tabular}
\end{center}

\section*{An evidence hierarchy}

Not every measurement below carries the same weight, and the distinction is used
throughout rather than restated each time. \S\ref{sec:synthesis} develops it in full.

\begin{description}
\item[State-level quantities] --- the energy, the local-energy variance $\Var(E_L)$, and the
  overlap between two trained states. These are properties of the wavefunction and are
  unaffected by how the parameters are arranged. The architectural conclusions rest here.
\item[Controlled interventions] --- deleting a component and pricing the loss in energy, or
  swapping one architecture for another and retraining. Reproducible for a fixed model and
  protocol, but conditional on how that model divided the correlation between its parts.
\item[Component-level diagnostics] --- ranks, displacement magnitudes, per-branch loadings.
  These are interpretive and \emph{gauge-like}: the Jastrow and the backflow are partial
  substitutes, so the division of labour between them records the training route as much as
  the state. Reported where they illuminate a mechanism; never load-bearing.
\end{description}

\section*{Three different notions of ``dimension''}

Three distinct geometries are measured in these chapters and all three are conventionally
described as low- or high-dimensional. They decompose different objects and are not
comparable with one another.

\begin{center}\small
\setlength{\tabcolsep}{5pt}
\begin{tabular}{@{}p{0.20\textwidth} p{0.24\textwidth} p{0.46\textwidth}@{}}
\toprule
Quantity & Object decomposed & What it counts \\
\midrule
Latent rank $r_{\mathrm{eff}}(Z)$ & correlator feature covariance &
  feature-space directions the correlator uses \\
Displacement rank & backflow displacement covariance &
  spatial directions the backflow moves particles along \\
Tangent dimension $d_{\mathrm{eff}}(S)$ & quantum geometric tensor &
  parameter-space directions carrying Fisher weight \\
\bottomrule
\end{tabular}
\end{center}

\noindent A correlator that is low-rank in its \emph{features} says nothing about how many
\emph{tangent} directions it uses, and neither constrains how many \emph{spatial} directions
the backflow occupies. Where the text says ``rank'' without qualification it means the one
named in the surrounding table. Two spectral counts are also in use and differ: the entropy
effective rank $\exp(H_1)$ and the participation ratio $(\sum_a\lambda_a)^2/\sum_a\lambda_a^2$,
with the second never larger than the first (\S\ref{subsec:repr-erank}).

\section*{Units and ordering}

Structural quantities are accumulated in trap units $\tilde{\mathbf r}=\sqrt\omega\,\mathbf r$
so that shapes can be compared across confinements; absolute radii quoted against classical
predictions are in Bohr; and sampling-proposal geometry is expressed in oscillator lengths
$\ell=\omega^{-1/2}$. Each table states which. Tables and prose run from strong to weak
confinement, $\omega:1\to0.5\to0.1\to0.01\to10^{-3}$, which is the direction the physics
develops.

\section*{Naming}

In these chapters the two learnable objects are named by function rather than by
implementation class. The \emph{correlator} is the Jastrow factor $W_\theta$ and comes in a
\emph{separable} form (independent per-particle and pairwise embeddings, then pooling) and a
\emph{message-passing} form. The \emph{backflow} is the coordinate displacement field and
comes in a \emph{conventional} form (a single round of pairwise messages) and a
\emph{message-passing} form (persistent, transported edge state). The corresponding code
classes---\texttt{CTNNJastrowVCycle}, \texttt{BackflowNet}, \texttt{CTNNBackflowNet}---are
specified in Chapter~\ref{ch:trial_wavefunction} and used here only where the exact
implementation matters.

\clearpage
